\documentclass{article}
\usepackage{decomp}
\begin{document}

\decompnote{Full verbatim transcript of US patent 7,146,895 B2, produced tool-first: the text was placed from a Datalab API (accurate mode) OCR of the USPTO image scan made on 2026-10-20 (\texttt{work/datalab/US7146895B2.md}) by the script \texttt{work/build\_main.py}, and proofread against the page images by the decompiler, with the Google Patents text of the same patent (\texttt{source/US7146895B2.google\_patents.txt}) as a second opinion for doubtful words; every correction is in the script's correction table and in the README. The source PDF is the USPTO full-text image: 8 pages, one 300-ppi bilevel scan per page, no text layer; the page images are authoritative. The patent is a public USPTO document and prints no copyright notice. Left out as layout: the running header of the drawing sheets (``U.S. Patent \quad Dec. 12, 2006 \quad Sheet $n$ of 4 \quad US 7,146,895 B2''), the running header of the text pages (``US 7,146,895 B2''), the column numbers 1--6 printed at the head of each text column, and the line numbers (5, 10, 15, \ldots) printed in the gutter between the columns.}

\sourcefile{US7146895B2.pdf}

\pdfpage{1}
\noindent\texttt{US007146895B2}
\decompnote{The line above is printed under the barcode at the top of the front page; the OCR dropped it and the decompiler supplied it from the page image. The front page is set in two columns; below, the bibliographic fields follow the left column, the U.S. patent documents list (which runs from the foot of the left column into the right column) is given as one table, and the examiner lines and the abstract follow from the right column.}

\noindent (12) {\large\textbf{United States Patent}}\\
\hspace*{2.2em}\textbf{Kong et al.}

\noindent (10) \textbf{Patent No.:} \textbf{US 7,146,895 B2}\\
(45) \textbf{Date of Patent:} \textbf{Dec. 12, 2006}

\noindent (54) \textbf{SLIDING BLADE MICROTOME}

\noindent (76) Inventors: \textbf{George Y. Kong}, 116 Redonovan Dr., Santa Clara, CA (US) 95051; \textbf{Jian-Qiang Kong}, 985-C5 Spring Forest Rd., Greenville, NC (US) 27834

\noindent ( * ) Notice: Subject to any disclaimer, the term of this patent is extended or adjusted under 35 U.S.C. 154(b) by 201 days.

\noindent (21) Appl. No.: \textbf{10/969,350}

\noindent (22) Filed: \textbf{Oct. 21, 2004}

\noindent (65) \textbf{Prior Publication Data}\\
US 2006/0086221 A1 Apr. 27, 2006

\noindent (51) \textbf{Int. Cl.}\\
\textbf{\textit{B26D 1/10}} (2006.01)

\noindent (52) \textbf{U.S. Cl.} ..... \textbf{83/705}; 83/409.2; 83/703

\noindent (58) \textbf{Field of Classification Search} ..... 83/915.5, 83/703, 409.2, 705, 563, 613, 638 See application file for complete search history.

\noindent (56) \textbf{References Cited}

\noindent U.S. PATENT DOCUMENTS\par\nopagebreak
\noindent \begin{tabular}{@{}l@{\quad}l@{\quad}r@{\quad}l@{}}
4,455,910 & A & 6/1984 & Kraft et al.\\
4,625,608 & A & 12/1986 & Behme et al.\\
4,754,675 & A & 7/1988 & Segal\\
5,148,729 & A & 9/1992 & Krumdieck\\
5,161,446 & A & 11/1992 & Holbl et al.\\
5,461,953 & A & 10/1995 & McCormick\\
5,522,294 & A & 6/1996 & Krumdieck\\
5,533,342 & A & 7/1996 & Gordon\\
5,550,033 & A & 8/1996 & Krumdieck\\
5,713,255 & A & 2/1998 & Izvozichikov et al.\\
5,740,708 & A & 4/1998 & Tabone\\
6,209,437 & B1 & 4/2001 & Izvozichikov et al.\\
6,253,653 & B1 & 7/2001 & Walter et al.\\
6,598,507 & B1 & 7/2003 & Gunther et al.\\
6,634,268 & B1 & 10/2003 & Guenther et al.\\
6,651,538 & B1 & 11/2003 & Tamura et al.\\
2003/0101858 & A1 & 6/2003 & Tamura et al.\\
\end{tabular}

\noindent FOREIGN PATENT DOCUMENTS\par\nopagebreak
\noindent \begin{tabular}{@{}l@{\quad}l@{\quad}r@{\quad}l@{}}
DE & 31 27 266 & A1 & 1/1983\\
GB & 2 182 786 & A & 5/1987\\
\end{tabular}

\noindent \textit{Primary Examiner}—Charles Goodman\\
(74) \textit{Attorney, Agent, or Firm}—Richard C. Litman

\noindent (57) \textbf{ABSTRACT}

The sliding blade microtome is a tissue specimen cutting device including a specimen syringe and an adjustable cutting mechanism for diagonally slicing the specimen. The specimen syringe is a cylindrical shaft having a first end and a second end for holding and supporting a specimen. The first end has a diameter that is smaller than the diameter of the shaft, forming a compressed lip for applying pressure on the specimen prior to slicing by the blade. The pressure used to compress the gelatin and the specimen holds the gelatin and the specimen in a sturdy position to slice the specimen to precision without chatter marks. The specimen is pushed out of the first end of the specimen syringe by a motor mechanism or a manually adjustable micrometer drive.

\begin{center}
\textbf{11 Claims, 4 Drawing Sheets}
\end{center}

\begin{center}
\includegraphics[width=0.8\linewidth,height=0.5\textheight,keepaspectratio]{figures/front.png}
\end{center}
\alttext{Representative drawing on the front page (below ``11 Claims, 4 Drawing Sheets''): a black-and-white line drawing, the same sectional side view of the first embodiment as FIG. 3 on Sheet 3, printed upright and at a smaller scale, with the same 32 reference numerals in italic: 100 (the sliding blade microtome, arrow at the upper right), 115 (base), 120 (bath), 121 (buffer medium, the wavy liquid line in the bath), 122 (syringe aperture in the bath wall), 124 (sealant ring), 126 (syringe mount), 128 (spacer ring), 129 (stopper), 130 (specimen syringe), 132 (first end of the syringe, the compressed lip inside the bath), 134 (second end of the syringe), 140 (plunger), 142 (first end of the plunger), 144 (second end of the plunger), 146 (gelatin), 148 (specimen), 150 (cutting mechanism, arrow at the upper left), 152 (support arm), 153 (first end of the linear way guide), 154 (linear way guide, underlined), 156 (mounting accessory), 158 (cross clamp), 160 (axis rod), 162 (blade holder), 164 (screws), 166 (blade clamping plate), 168 (blade), 170 (cutting plane, the dashed vertical line at the blade tip), 172 (micrometer drive), 174 (micrometer mounting block), 176 (micrometer tip). See the description of FIG. 3 for what each part looks like and where it sits; the only differences are the scale and the orientation (this copy is upright, Sheet 3 is printed sideways). No caption or figure label is printed with it.}

\pdfpage{2}
\begin{center}
\rotatebox[origin=c]{-90}{\includegraphics[height=0.95\linewidth,width=0.72\textheight,keepaspectratio]{figures/fig1.png}}
\end{center}
\noindent\textit{FIG. 1}
\alttext{Sheet 1 of 4, FIG. 1: ``an environmental, perspective view of a sliding blade microtome'' (first embodiment, 100). Black-and-white line drawing printed sideways on the sheet (the label ``FIG. 1'' and all numerals read from bottom to top; the figure is upright when the sheet is turned a quarter turn clockwise, as the PNG is rotated here). Reference numerals in italic, each with a curved leader line to its part; underlined 154 is written on the part itself. Seen upright: the base 115 is a flat rectangular slab drawn in perspective, its long side running from the lower left to the right, with short strokes suggesting a surface finish. On its left half stands the bath 120, an open-top rectangular box whose faces carry diagonal strokes; the leader of 122 (syringe aperture) points to a round hole in the bath's right-hand wall, through which the short cylindrical first end 132 of the specimen syringe projects into the bath. Behind and to the right of the bath rises the support arm 152, a tall upright plate; bolted to its front face is the linear way guide 154, a long straight bar with a row of about eight round screw heads along its length, inclined so that its first end 153 (upper left, labelled at the top of the bar) is higher than its second end 155 (lower right, labelled where the bar ends just above the bath). On the upper end of the guide rides the mounting accessory 156 (a block), with the cross clamp 158 (a vertical plate with a screw head) hanging from it; the axis rod 160 is a long round rod projecting horizontally to the left from the cross clamp. Below the cross clamp the blade holder 162 (a small block with a thumbscrew) carries the blade clamping plate 166 (a plate with two holes) fixed by the screws 164 (two screw heads, one leader forked to both); the blade 168 is the thin strip projecting from the lower end of the plate towards the bath. Arrow 150 (upper left) points at this whole cutting mechanism. To the right of the bath a block with a thumbscrew on top (the syringe mount, not numbered in this view) holds the specimen syringe 130, a horizontal cylinder; its second end 134 is at the right. From there a thin rod runs to the right through an upright block with a round hole (the micrometer mounting block, not numbered in this view; it has a thumbscrew on its right side) to the micrometer drive 172, a micrometer with a knurled thimble at the far right. Arrow 100 (upper right) points at the whole microtome. ``FIG. 1'' is printed at the lower right. Not a plotted panel: no data file.}

\pdfpage{3}
\begin{center}
\rotatebox[origin=c]{-90}{\includegraphics[height=0.95\linewidth,width=0.72\textheight,keepaspectratio]{figures/fig2.png}}
\end{center}
\noindent\textit{FIG. 2}
\alttext{Sheet 2 of 4, FIG. 2: ``a front view of the sliding blade microtome'' (first embodiment, 100), line drawing printed sideways like Sheet 1 (upright when turned a quarter turn clockwise). Seen upright: the base 115 is the long horizontal bar across the bottom of the drawing. On it, left of centre, stands the support arm 152, a tall rectangle with diagonal strokes (its numeral is underlined on the part); the linear way guide 154 is a bar fastened across the front of the arm, sloping down from the upper left to the lower right, with a row of seven small circled screw heads along it and two plain circles near its upper end; its first end 153 is the upper left end and its second end 155 the lower right end, which reaches down to just above the bath. At the upper end of the guide sits the mounting accessory 156 (a short block); from it hangs the cross clamp 158, a vertical bar with a circle at its top and bottom, the top circle being the end of the axis rod 160 seen end-on. A pair of arrowheads drawn along the guide beside 160 and two short arrows beside 153 and 156 mark the directions of sliding and of tilting. Below the cross clamp is the blade holder 162, a rectangle holding the blade 168 (the horizontal strip along its lower edge) with the two screws 164 drawn as circled crosses. A vertical dash-dot section line, marked ``3'' with an arrow at its top end and at its bottom end (both arrows pointing left), passes through the blade holder and the bath: it is the section plane of FIG. 3. To the right stands the bath 120, a rectangle with diagonal strokes and a dashed inner outline (the inside of the walls); a broken-away patch with a wavy outline in its front wall shows the syringe end-on as a set of concentric circles: the syringe aperture 122 (the hole in the wall), the first end 132 of the specimen syringe 130 (the rings), and inside them the gelatin 146 with the specimen 148 at the centre, drawn as a small irregular shape; a small circle on each side of the rings is a screw hole. Above the bath is the top of the syringe mount with a thumbscrew (not numbered). Arrow 100 at the upper right points at the whole. ``FIG. 2'' is printed at the lower right. Not a plotted panel: no data file.}

\pdfpage{4}
\begin{center}
\rotatebox[origin=c]{-90}{\includegraphics[height=0.95\linewidth,width=0.72\textheight,keepaspectratio]{figures/fig3.png}}
\end{center}
\noindent\textit{FIG. 3}
\alttext{Sheet 3 of 4, FIG. 3: ``a sectional side view of the sliding blade microtome'' (first embodiment, 100; the section 3--3 marked in FIG. 2), line drawing printed sideways like Sheets 1 and 2 (upright when turned a quarter turn clockwise). Cut parts are hatched with diagonal lines. Seen upright: the base 115 is a hatched slab across the bottom. At the left stands the bath 120, a U-shaped hatched vessel in section; the buffer medium 121 is shown by a wavy line (the liquid surface) across the upper part of the bath. The specimen syringe 130 is a horizontal tube entering the bath from the right through the syringe aperture 122 (leader from below, pointing at the hole in the bath's right wall); its first end 132, inside the bath, is drawn with the wall turned inwards to a narrower opening (the compressed lip), and its second end 134 is at the right. Inside the syringe lies the plunger 140, a hatched rod, with its first end 142 at the left and its second end 144 at the right; at the lip, the gelatin 146 (stippled) with the specimen 148 (the stippled block) is being pushed out of the first end, its left face just beyond the lip. Where the syringe passes through the bath wall, the sealant ring 124 is drawn as a small circle above and below the tube (an O-ring in section); to the right of the bath wall the hatched syringe mount 126 surrounds the tube; the spacer ring 128 and the stopper 129 are the small stepped pieces at the right-hand face of the mount, on the tube. Further right the micrometer mounting block 174, an upright block with a round hole standing on the raised part of the base and carrying a thumbscrew on its right side, holds the micrometer drive 172 (a micrometer with a knurled thimble at the far right); its thin micrometer tip 176 reaches to the left and touches the second end 144 of the plunger. Above the syringe mount rises the support arm 152, a tall hatched column; the linear way guide 154 (underlined, seen end-on as a rectangle at the top, with its first end 153 labelled at the top) sits against the arm, and to its left the mounting accessory 156 (a stepped block) and the cross clamp 158 (a vertical plate) hang down; the axis rod 160 is a horizontal bar projecting to the left from the cross clamp, with a screw head where they meet. Below, the blade holder 162 (a block with a thumbscrew on its left) holds the blade 168, a hatched strip angled down and to the right, together with the blade clamping plate 166 and the screws 164 (one cross-hatched screw drawn across the blade). The cutting plane 170 is a dashed vertical line through the blade tip and the lip of the syringe: the tip of the blade sits on it, so the blade makes an acute angle with the plane. Arrow 150 (upper left) points at the cutting mechanism, arrow 100 (upper right) at the whole. ``FIG. 3'' is printed at the lower right. The same drawing, upright and smaller, is the representative figure on the front page. Not a plotted panel: no data file.}

\pdfpage{5}
\begin{center}
\rotatebox[origin=c]{-90}{\includegraphics[height=0.95\linewidth,width=0.72\textheight,keepaspectratio]{figures/fig4.png}}
\end{center}
\noindent\textit{FIG. 4}
\alttext{Sheet 4 of 4, FIG. 4: ``a sectional side view of a sliding blade microtome with a mechanical control mechanism'' (second embodiment, the automated microtome 200), a schematic line drawing printed sideways like the other sheets (upright when turned a quarter turn clockwise). Seen upright: a base (not numbered) across the bottom, with a raised block at its right half. On the left half stands the bath (not numbered), a rectangle with diagonal strokes, a dashed inner outline and a wavy liquid line; the specimen syringe (not numbered) enters it through the right wall, its narrowed lip drawn inside the bath, and is held by a block on the raised part of the base. Behind the syringe the plunger 236 is a horizontal cylinder (its numeral underlined on the part); 234 is a smaller collar at the plunger's right end (this numeral is printed in the drawing but is not named in the description). Below and to the right, the lead screw block 254 (a rectangle, numeral underlined) rides on the lead screw 256, a long horizontal toothed bar running to the motor 252 (a box at the far right, numeral underlined); a short arrow under 254 points left, the direction in which the plunger is advanced. The sensor 250 is a small rectangle hanging above the lead screw block on a line from the control unit. The control unit 220 is a box at the upper right (numeral underlined) with four lines (wires) leaving it: down to the motor 252, down-left to the sensor 250, and across the top of the drawing to the left to the sensor 240 and the motor 242. At the upper left, the cutting mechanism 232 is drawn in dashed (phantom) outline as a stepped shape with an arrow from 232 pointing at it; 230, with a dashed leader, is a small dashed outline at the lower left of that shape (printed but not named in the description). The blade (not numbered) is drawn in solid lines below the phantom outline, angled down and to the right, its tip on the cutting plane 170, a dashed vertical line reaching down into the bath. The motor 242 (a box, numeral underlined) sits above the blade; from it the lead screw 246, a vertical toothed bar, runs down to the lead screw block 244 (a box, numeral underlined), with a short arrow under 244 pointing down, the direction of the cut; the sensor 240 is a small rectangle just to the right of 244 on its line from the control unit. Arrow 200 (upper right) points at the whole. ``FIG. 4'' is printed at the lower right. Not a plotted panel: no data file.}

\pdfpage{6}
\decompnote{PDF pages 6--8 are the text of the patent in two columns (columns 1--6). The column numbers and the line numbers (every fifth line, printed in the gutter) are left out; paragraphs broken by a column or page break are joined, with the \pdfpage{} marker inline. Reference numerals are printed in bold, as kept here.}
\section*{SLIDING BLADE MICROTOME}

\subsection*{BACKGROUND OF THE INVENTION}

\subsubsection*{1. Field of the Invention}

The present invention relates to microtomes, more specifically to a sliding blade microtome.

\subsubsection*{2. Description of the Related Art}

Much modern scientific research involves microscopic analysis. Much of today's microscopic analysis requires tissues to be prepared by devices that provide thin specimens, such as microtomes. Many microtomes are specifically designed for preparing specimen for certain types of microscopes. For instance, rotary microtomes cut thin sections of specimen for light microscopy and ultramicrotomes are used to prepare ultrathin sections for light and electron microscopy. Live tissues, such as brain or liver tissues, are prepared by oscillating or vibrating blade microtomes called vibratomes. Vibratomes are currently employed in biomedical studies such as electrophysiology recording, organotypic tissue culture and immunohistochemistry experiments. Though vibratomes have advanced the study of live and fresh tissues, there are disadvantages, such as slow cutting speeds, inconsistent slice thickness, lack of a support for the specimen and the presence of chatter marks on the sliced specimen, that leave the vibratome lacking as a fully efficient microtome. A microtome is desired that will uniformly slice specimen at a specific thickness without imparting chatter marks on the slices. The microtome should not only securely hold, but also steadily support the specimen. Thus, a microtome solving the aforementioned problems is desired.

\subsection*{SUMMARY OF THE INVENTION}

The microtome is a tissue specimen-cutting device for producing high quality specimen slices. The device includes a bath, a specimen syringe for holding, supporting and compressing the specimen, and an adjustable cutting mechanism for diagonally slicing the specimen. The cutting mechanism has a linear way guide, a mounting accessory slidingly attached to the linear way guide, and a cross clamp attached to the mounting accessory, a blade holder being adjustably attached to the cross clamp and holding a blade.

The specimen syringe is a cylindrical shaft having a first end and a second end for holding and supporting a specimen. The shaft and the second end have the same diameter. The first end has a diameter that is smaller than the diameter of the shaft and the second end, forming a compressed lip at the first end. A specimen is prepared by gluing it to a first end of a plunger and embedding the specimen in a gelatin, such as low gel point agarose or similar material. The specimen is pushed out of the first end of the specimen syringe by placing pressure on a second end of the plunger using a motor mechanism or a manually adjustable micrometer drive.

Thus, as the specimen is pushed out of the compressed lip at the first end of the specimen syringe, pressure is applied to the specimen and the gelatin. This pressure functions to compress the gelatin and the specimen when they are forced out of the opening of the specimen syringe, thereby clamping and holding the gelatin and the specimen in a sturdy position to slice the specimen to precision without chatter marks.

The blade of the cutting assembly is disposed adjacent the first end of the specimen syringe. A cutting plane is defined along the first end of the specimen syringe. The blade can be angled by rotating the blade holder about the cross clamp to define a cutting angle against the cutting plane.

Ideally the first end and a portion of the syringe shaft are disposed in a bath through a syringe aperture. The syringe aperture is completely submerged under a buffer medium, which is held in the bath. A sealant ring seals the gap between the specimen syringe and the syringe aperture of the bath, thereby preventing the buffer medium from leaking.

A spacer ring surrounds the specimen syringe and works together with a stopper, which is disposed on the syringe, to determine how much of the syringe will be exposed in the bath. The syringe slides into the bath until the stopper abuts the spacer ring and impedes the syringe from moving. The width of the spacer ring, therefore, directly determines how much of the syringe is exposed. A wide spacer ring allows only a small portion of the syringe to extend into the bath, and a narrow spacer ring allows a larger portion of the syringe to extend into the bath. The specimen, as well as the tissue slices cut by the microtome, is disposed in the bath, where the buffer medium not only provides oxygen and nutrition to the specimen, but also lubricates the blade as it cuts the specimen in the bath.

These and other features of the present invention will become readily apparent upon consideration of the following specification and drawings.

\subsection*{BRIEF DESCRIPTION OF THE DRAWINGS}

FIG. \textbf{1} is an environmental, perspective view of a sliding blade microtome according to the present invention.

FIG. \textbf{2} is a front view of the sliding blade microtome according to the present invention.

FIG. \textbf{3} is a sectional side view of the sliding blade microtome according to the present invention.

FIG. \textbf{4} is a sectional side view of a sliding blade microtome with a mechanical control mechanism according to the present invention.

Similar reference characters denote corresponding features consistently throughout the attached drawings.

\subsection*{DETAILED DESCRIPTION OF THE PREFERRED EMBODIMENTS}

The present invention is a sliding blade microtome, a first embodiment of which is designated as \textbf{100} in the figures. The sliding blade microtome \textbf{100} is a tissue specimen-cutting device for slicing live and pre-fixed tissue. The microtome \textbf{100} has a base \textbf{115} for supporting a specimen syringe \textbf{130} having a first end \textbf{132} and a second end \textbf{134}, a bath \textbf{120} having a syringe aperture \textbf{122}, an adjustable cutting mechanism \textbf{150} and a micrometer drive \textbf{172}. The first end \textbf{132} of the specimen syringe \textbf{130} is ideally disposed through the syringe aperture \textbf{122} of the bath \textbf{120}. The second end \textbf{134} of the syringe \textbf{130} is disposed adjacent the micrometer drive \textbf{172}.

Referring to FIGS. \textbf{1}–\textbf{3}, the cutting mechanism \textbf{150} is shown comprising a mounting accessory \textbf{156} having an axis rod \textbf{160}, a cross clamp \textbf{158}, a blade holder \textbf{162} and a blade \textbf{168}. The cross clamp \textbf{158} is disposed on the axis rod \textbf{160} of the mounting accessory \textbf{156}. The cross clamp \textbf{158} serves as a support to which the blade holder \textbf{162} is pivotally attached. The blade \textbf{168} is clamped to the blade holder \textbf{162} by a blade clamping plate \textbf{166} and is tightened in position using screws \textbf{164}.

The cutting mechanism \textbf{150} is disposed on a linear way guide \textbf{154} that is mounted on a support arm \textbf{152} of the base \pdfpage{7} \textbf{115}. The mounting accessory \textbf{156} is the portion of the cutting mechanism \textbf{150} that directly contacts the guide \textbf{154}. Specifically, the mounting accessory \textbf{156} allows the cutting mechanism \textbf{150} to slide from a first end \textbf{153} of the linear way guide \textbf{154} to a second end \textbf{155}. The second end \textbf{155} of the linear way guide \textbf{154} is disposed adjacent the first end \textbf{132} of the specimen syringe \textbf{130} that is situated within the bath \textbf{120}. The linear way guide \textbf{154} may be tilted up or down to dispose the guide \textbf{154} at any angle above the bath \textbf{120}. The linear way guide \textbf{154} is pivoted about a bolt that is secured through one of a number of holes in the support arm \textbf{152}. The first end \textbf{153} of the guide \textbf{154} may be level with the second end \textbf{155}, or may be raised above the second end \textbf{155}.

Referring now to FIG. \textbf{3}, the specimen syringe \textbf{130} is a cylindrical, hollow shaft designed for supporting and holding a specimen \textbf{148}. The specimen syringe \textbf{130} is positioned and clamped in place on an immovable syringe mount \textbf{126} on the base \textbf{115}. The first end \textbf{132} of the syringe \textbf{130} has a diameter that is smaller than the diameter of the second end \textbf{134} and the body of the syringe \textbf{130}, and consequentially defines a compressed lip. The compressed lip at the first end \textbf{132} has a smaller cross section than the rest of the syringe \textbf{130}.

The specimen syringe \textbf{130} holds the specimen \textbf{148} after it is prepared on a plunger \textbf{140}. The specimen \textbf{148} is prepared by being fixed to a first end \textbf{142} of the plunger \textbf{140} with an adhesive, such as isocyanate glue or the like. Afterwards, the specimen \textbf{148} is embedded in a gelatin \textbf{146}, such as a low gel point agarose, molten agarose or similar material once the specimen \textbf{148} is affixed to the plunger \textbf{140} in the syringe \textbf{130}.

The compressed lip at the first end \textbf{132} of the syringe \textbf{130} is the mechanism through which pressure is placed on the gelatin \textbf{146} and the specimen \textbf{148} as they exit the syringe \textbf{130}. The pressure applied by the compressed lip restricts movement of the specimen \textbf{148} and the gelatin \textbf{146}, holds both the gelatin \textbf{146} and the specimen \textbf{148} in a firm, sturdy position, and mechanically solidifies the two into a firmer mass. The importance of applying pressure to the gelatin \textbf{146} and the specimen \textbf{148} prior to being cut is to counteract mechanical stresses of the blade \textbf{168} as it cuts the specimen \textbf{148}. Mechanical stress can cause tissue distortion and cell damage of the specimen \textbf{148} during slicing. Thus, due to the pressure from the compressed lip of the first end \textbf{132}, the specimen \textbf{148} can be precisely sliced without chatter marks.

Once the gelatin \textbf{146} and the specimen \textbf{148} are pushed past the compressed lip of the first end \textbf{132} of the syringe \textbf{130}, the operator can slide the blade \textbf{168} from the first end \textbf{153} to the second end \textbf{155} of the linear way guide \textbf{154} to cut the specimen \textbf{148}. The angle of the linear way guide \textbf{154} over the first end \textbf{132} of the syringe \textbf{130} determines how the specimen \textbf{148} is cut. If the guide \textbf{154} is angled so that the first end \textbf{153} is disposed higher than the second end \textbf{155}, then the blade \textbf{168} cuts the specimen \textbf{148} diagonally, since the angle of the guide \textbf{154} relative to the support \textbf{152} is locked while sliding the cutting mechanism \textbf{150} along the guide \textbf{154}.

The blade \textbf{168} slides along the guide \textbf{154}, abutting the first end \textbf{132} of the syringe \textbf{130}. Specifically, the blade \textbf{168}, via the blade holder \textbf{162}, is rotated at an angle from a cutting plane \textbf{170}, which is the outermost end of the first end \textbf{132}. The blade \textbf{168} forms an acute angel with the cutting plane \textbf{170}. The cutting plane \textbf{170} is perpendicular with the axially oriented syringe \textbf{130}. The cutting angle may be adjusted by rotating the blade holder \textbf{162} in the cross clamp \textbf{158} to increase or decrease the angle between the blade \textbf{168} and the cutting plane \textbf{170}. Once the blade \textbf{168} cuts slices from the specimen \textbf{148}, they fall into the bath \textbf{120}.

As described above, the first end \textbf{132} and a portion of the shaft of the syringe \textbf{130} are disposed through the syringe aperture \textbf{122} of the bath \textbf{120}. The bath \textbf{120} has a bottom wall and at least one peripheral wall extending around the bottom wall; the bath holds a buffer medium \textbf{121} such as saline. The syringe aperture \textbf{122} is completely submerged under the buffer medium \textbf{121}. To prevent the buffer medium \textbf{121} from leaking a sealant ring \textbf{124} seals the gap between the specimen syringe \textbf{130} and the wall of the bath \textbf{120}. The sealant ring \textbf{124} is disposed adjacent the syringe aperture \textbf{122} and is sandwiched between the bath \textbf{120} and the syringe mount \textbf{126}. The buffer medium \textbf{121} both lubricates the blade \textbf{168} as it cuts the specimen \textbf{148} and provides oxygen and nutrition to the specimen \textbf{148} and the specimen slices cut by the microtome \textbf{100}. In alternative embodiments, the bath \textbf{120} may be disposed below the syringe \textbf{130} to catch the specimen slices that are cut by the microtome \textbf{100}.

The portion of the syringe \textbf{130} that is disposed through the syringe aperture \textbf{122} is determined by a stopper \textbf{129} that is welded or otherwise disposed on the outside of the syringe \textbf{130}, and by a spacer ring \textbf{128}. The spacer ring \textbf{128} surrounds the specimen syringe \textbf{130} and is sandwiched between the syringe mount \textbf{126} and the stopper \textbf{129}. The syringe \textbf{130} moves axially through the syringe aperture \textbf{122} until the stopper \textbf{129} abuts the spacer ring \textbf{128} and stops the syringe \textbf{130}. The width of the spacer ring \textbf{128} is a significant factor in determining how much of the syringe \textbf{130} is exposed in the bath \textbf{120}. For instance, a wide spacer ring \textbf{128} allows only a small portion of the syringe \textbf{130} to extend into the bath \textbf{120}. Narrow spacer rings \textbf{128} allow a larger portion of the syringe \textbf{130} to extend into the bath \textbf{120}.

The thickness of the slices cut from the specimen \textbf{148} may be controlled by the axial movement of the plunger \textbf{140} through the specimen syringe \textbf{130}. The axial movement of the plunger \textbf{140} through the specimen syringe \textbf{130} may be controlled manually or mechanically. As shown in FIGS. \textbf{1}–\textbf{3}, the specimen \textbf{148} may be pushed out of the first end \textbf{132} of the specimen syringe \textbf{130} using a micrometer drive \textbf{172}. A micrometer tip \textbf{176} of the micrometer drive \textbf{172} is positioned at a second end \textbf{144} of the plunger \textbf{140} and incrementally pushes the specimen \textbf{148} out of the first end \textbf{132} of the syringe \textbf{130}. The micrometer drive \textbf{172} is mounted on a micrometer mounting block \textbf{174} of the base \textbf{115}.

FIG. \textbf{4} shows a second embodiment of the microtome, which is an automated microtome \textbf{200} that drives the plunger \textbf{236} through the syringe \textbf{130} via a control unit \textbf{220}, sensors \textbf{240}, \textbf{250}, motors \textbf{242}, \textbf{252}, lead screw blocks \textbf{244}, \textbf{254} and lead screw \textbf{246}, \textbf{256}. The motor \textbf{242} drives cutting mechanism \textbf{232} through lead screw \textbf{246} and lead screw block \textbf{244}. The motor \textbf{252} operates the plunger \textbf{236}, moving it forward using lead screw \textbf{256} and lead screw block \textbf{254}. The control unit \textbf{220} monitors the positions of the lead screw blocks \textbf{244}, \textbf{254} through sensors \textbf{240}, \textbf{250}, respectively, and sends commands to control the motors \textbf{242}, \textbf{252}, respectively.

The microtomes \textbf{100}, \textbf{200} differ in how the plungers \textbf{140} and \textbf{236}, respectively, are pushed out of the syringe \textbf{130}. However, both microtomes \textbf{100}, \textbf{200} hold and support the specimen \textbf{148} using the specimen syringe \textbf{130}. Also, both have the compressed lip at the first end \textbf{132} to compress the specimen \textbf{148} prior to being sliced by the blade \textbf{168}. Furthermore, the cutting mechanism \textbf{150}, \textbf{232} of each microtome \textbf{100}, \textbf{200} is directed to slidingly cut specimen \textbf{148} \pdfpage{8} directly within the bath \textbf{120} as the blade \textbf{168} slides from the first end \textbf{153} to the second end \textbf{155} of the linear way guide \textbf{154}.

It is to be understood that the present invention is not limited to the embodiments described above, but encompasses any and all embodiments within the scope of the following claims.

We claim:

\begin{hangpars}
\textbf{1}. A microtome, comprising:

{\leftskip=1.5em\parindent=0pt a base;\par}
{\leftskip=1.5em\parindent=0pt a bath disposed on the base, the bath being adapted for holding a buffer medium, the bath having a bottom wall and at least one peripheral wall extending around the bottom wall, the bath having a syringe aperture defined therein;\par}
{\leftskip=1.5em\parindent=0pt a specimen syringe having a hollow shaft extending through the syringe aperture of the bath, wherein the shaft of the syringe has a first end and a second end, the first end extending into the bath, the first end having a constricted diameter forming a lip for compressing the specimen exiting the first end of the syringe;\par}
{\leftskip=1.5em\parindent=0pt a plunger having a first end and a second end, the first end of the plunger forming a holding surface adapted for having a specimen secured thereto, the plunger being slidably disposed in the specimen syringe; and\par}
{\leftskip=1.5em\parindent=0pt a cutting mechanism attached to the base adjacent the first end of the specimen syringe;\par}
{\leftskip=1.5em\parindent=0pt whereby the specimen is pushed through the syringe shaft into the bath and the cutting mechanism slices the specimen within the bath.\par}

\textbf{2}. The microtome of claim \textbf{1}, wherein the specimen syringe is defined in the peripheral wall of said bath.

\textbf{3}. The microtome of claim \textbf{1}, wherein the cutting mechanism comprises:

{\leftskip=1.5em\parindent=0pt a linear way guide attached to the base;\par}
{\leftskip=1.5em\parindent=0pt a mounting accessory slidingly attached to the linear way guide;\par}
{\leftskip=1.5em\parindent=0pt a cross clamp attached to the mounting accessory;\par}
{\leftskip=1.5em\parindent=0pt a blade holder attached to the cross clamp; and\par}
{\leftskip=1.5em\parindent=0pt a blade attached to the blade holder.\par}

\textbf{4}. The microtome of claim \textbf{3}, wherein the linear way guide is pivotally attached to the base.

\textbf{5}. The microtome of claim \textbf{3}, wherein the blade holder is pivotally attached to the cross clamp.

\textbf{6}. The microtome of claim \textbf{3}, wherein a cutting plane is defined normal to the first end of the shaft, said blade forming an acute angle with the cutting plane.

\textbf{7}. The microtome of claim \textbf{3}, wherein the mounting accessory is slidable along the linear way guide to a position in which said blade abuts the first end of the syringe shaft, thereby slicing the specimen into the bath.

\textbf{8}. The microtome of claim \textbf{1}, further comprising a micrometer drive attached to the base and abutting the second end of the plunger, the micrometer drive impelling the specimen out of the syringe into the bath.

\textbf{9}. A microtome comprising:

{\leftskip=1.5em\parindent=0pt a base;\par}
{\leftskip=1.5em\parindent=0pt a bath disposed on the base, the bath being adapted for holding a buffer medium, the bath having a bottom wall and at least one peripheral wall extending around the bottom wall, the bath having a syringe aperture defined therein;\par}
{\leftskip=1.5em\parindent=0pt a specimen syringe having a hollow shaft extending through the syringe aperture of the bath;\par}
{\leftskip=1.5em\parindent=0pt a plunger having a first end and a second end, the first end of the plunger forming a holding surface adapted for having a specimen secured thereto, the plunger being slidably disposed in the specimen syringe;\par}
{\leftskip=1.5em\parindent=0pt a cutting mechanism attached to the base adjacent the first end of the specimen syringe; and\par}
{\leftskip=1.5em\parindent=0pt a control unit comprising:\par}
{\leftskip=3.0em\parindent=0pt motorized means for sliding the cutting mechanism; and\par}
{\leftskip=3.0em\parindent=0pt motorized means for advancing the plunger in the syringe;\par}
{\leftskip=1.5em\parindent=0pt whereby the specimen is rushed through the syringe shaft into the bath and the cutting mechanism slices the specimen within the bath.\par}

\textbf{10}. A microtome, comprising:

{\leftskip=1.5em\parindent=0pt a base;\par}
{\leftskip=1.5em\parindent=0pt a specimen syringe mounted on the base, the specimen syringe being a hollow shaft having a first end and a second end, the first end having a compressed lip;\par}
{\leftskip=1.5em\parindent=0pt a bath disposed on the base, the bath being adapted for holding a buffer medium, the bath having a bottom wall and at least one peripheral wall extending around the bottom wall, the bath having a syringe aperture defined therein for slidably receiving the specimen syringe;\par}
{\leftskip=1.5em\parindent=0pt a plunger having a first end and a second end, the first end of the plunger being adapted for providing a holding surface to which a specimen is disposed; and\par}
{\leftskip=1.5em\parindent=0pt a cutting mechanism attached to the base and slidable across the first end of the specimen syringe.\par}

\textbf{11}. The microtome of claim \textbf{10}, wherein the cutting mechanism comprises:

{\leftskip=1.5em\parindent=0pt a linear way guide pivotally attached to the base;\par}
{\leftskip=1.5em\parindent=0pt a mounting accessory slidably attached to the linear way guide;\par}
{\leftskip=1.5em\parindent=0pt a cross clamp attached to the mounting accessory;\par}
{\leftskip=1.5em\parindent=0pt a blade holder adjustably attached to the cross clamp; and\par}
{\leftskip=1.5em\parindent=0pt a blade attached to the blade holder aligned with the compressed lip of the specimen syringe.\par}

\end{hangpars}

\begin{center}
* * * * *
\end{center}

\end{document}
